\documentclass[11pt]{article}
% Shared print layout. The web's provisional design and MathJax are independent.
\usepackage[utf8]{inputenc}
\usepackage[T1]{fontenc}
\usepackage[spanish,es-nodecimaldot]{babel}
\usepackage{lmodern}
\usepackage{amsmath,amssymb}
\usepackage{graphicx}
\usepackage{geometry}
\geometry{a4paper,left=20mm,right=20mm,top=25mm,bottom=22mm,headheight=16pt,headsep=8mm}
\usepackage{fancyhdr,lastpage}
\usepackage{enumitem}
\usepackage{xcolor}
\usepackage{pgfplots}
\pgfplotsset{compat=1.18}
\usepgfplotslibrary{fillbetween}
\usepackage[hidelinks,pdfencoding=auto]{hyperref}
\hypersetup{pdflang=es-ES,pdfauthor={},pdfsubject={Matemáticas Aplicadas a las Ciencias Sociales II}}
\pagestyle{fancy}
\fancyhf{}
\fancyhead[L]{\small Matemáticas Aplicadas a las Ciencias Sociales II}
\fancyhead[R]{\small 2.\textsuperscript{o} Bachillerato}
\fancyfoot[L]{\footnotesize Sistemas de ecuaciones I | Edición revisada 1.0}
\fancyfoot[R]{\footnotesize Página \thepage\ de \pageref{LastPage}}
\renewcommand{\headrulewidth}{0.4pt}
\setlength{\parindent}{0pt}
\setlength{\parskip}{0.5em}
\setlist{itemsep=0.3em,topsep=0.3em}
\newcommand{\apartado}[1]{\subsection*{#1}}
\pgfplotsset{
  revision/.style={
    width=0.94\linewidth,height=7.5cm,
    axis lines=middle,grid=major,grid style={gray!18},
    tick label style={font=\small},label style={font=\small},
    xlabel={$x$},ylabel={$y$},samples=160,
    legend style={at={(0.5,-0.16)},anchor=north,draw=none,font=\small},
    clip=true
  }
}

\hypersetup{pdftitle={Sistemas de ecuaciones I - Soluciones revisadas}}
\begin{document}
\section*{Sistemas de ecuaciones I: soluciones}
\section*{Ejercicio 1: sustitución regresiva}
La ecuación matricial equivale al sistema triangular
\[
\begin{cases}x+2y+3z=7,\\4y+5z=8,\\z=2.\end{cases}
\]
Resolvemos desde la última ecuación hacia la primera.
\subsection*{1. Tercera ecuación}
Directamente, $z=2$.
\subsection*{2. Segunda ecuación}
\[4y+5\cdot2=8\quad\Longrightarrow\quad4y=-2\quad\Longrightarrow\quad y=-\frac12.\]
\subsection*{3. Primera ecuación}
\[x+2\left(-\frac12\right)+3\cdot2=7\quad\Longrightarrow\quad x+5=7\quad\Longrightarrow\quad x=2.\]
Por tanto,
\[\boxed{(x,y,z)=\left(2,-\frac12,2\right)}.\]
\textbf{Comprobación completa:}
\[
\begin{pmatrix}1&2&3\\0&4&5\\0&0&1\end{pmatrix}
\begin{pmatrix}2\\-1/2\\2\end{pmatrix}
=\begin{pmatrix}2-1+6\\-2+10\\2\end{pmatrix}
=\begin{pmatrix}7\\8\\2\end{pmatrix}.
\]
La solución es única: la matriz triangular tiene determinante $1\cdot4\cdot1=4\neq0$.
\clearpage
\section*{Ejercicio 2: las tabletas}
Sean $A$, $B$ y $C$ las cantidades vendidas de cada tipo. Deben ser enteros no negativos.

\textbf{Interpretación del modelo.} El original llama ``coste por unidad'' a los importes de 120, 90 y 40 euros. Para relacionarlos con la recaudación, los interpretamos como los precios unitarios de venta del problema; no se incluyen márgenes ni otros ingresos.
\[
\begin{cases}
A+B+C=400,\\120A+90B+40C=36\,400,\\A+B=3C.
\end{cases}
\]
\subsection*{1. Obtener el número de tipo C}
De $A+B=3C$ y el total:
\[3C+C=400\quad\Longrightarrow\quad C=100,\qquad A+B=300.\]
\subsection*{2. Separar los tipos A y B}
Sustituyendo $C=100$ en la recaudación:
\[120A+90B=32\,400.\]
Con $B=300-A$,
\[120A+90(300-A)=32\,400\quad\Longrightarrow\quad30A=5\,400.\]
Así, $A=180$ y $B=120$.
\[\boxed{A=180,\qquad B=120,\qquad C=100}.\]
\textbf{Comprobación de las tres condiciones:}
\begin{align*}
180+120+100&=400,\\
120\cdot180+90\cdot120+40\cdot100
&=21\,600+10\,800+4\,000=36\,400,\\
180+120&=300=3\cdot100.
\end{align*}
Todos los valores son enteros y no negativos, coherentes con cantidades de tabletas.
\clearpage
\section*{Ejercicio 3(a): discusión según $\lambda$}
\[
A=\begin{pmatrix}1&1&1\\2&\lambda&3\\3&2&4\end{pmatrix},\qquad
b=\begin{pmatrix}2\\5\\\lambda\end{pmatrix}.
\]
El determinante es
\begin{align*}
\det A&=(4\lambda-6)-(8-9)+(4-3\lambda)\\
&=\lambda-1.
\end{align*}
\subsection*{Caso 1: $\lambda\neq1$}
Como $\det A\neq0$, $\operatorname{rg}(A)=\operatorname{rg}(A|b)=3$. El sistema es \textbf{compatible determinado}: tiene solución única.
\subsection*{Caso 2: $\lambda=1$}
No basta con que el determinante sea cero para decidir la compatibilidad. Aplicamos $F_2\leftarrow F_2-2F_1$ y $F_3\leftarrow F_3-3F_1$ a la matriz ampliada:
\[
\left(\begin{array}{ccc|c}1&1&1&2\\2&1&3&5\\3&2&4&1\end{array}\right)
\ \longrightarrow\ 
\left(\begin{array}{ccc|c}1&1&1&2\\0&-1&1&1\\0&-1&1&-5\end{array}\right).
\]
Ahora $F_3\leftarrow F_3-F_2$ produce
\[
\left(\begin{array}{ccc|c}1&1&1&2\\0&-1&1&1\\0&0&0&-6\end{array}\right).
\]
La última fila exige $0=-6$, una contradicción. Por tanto,
\[\operatorname{rg}(A)=2<3=\operatorname{rg}(A|b).\]
El sistema es \textbf{incompatible}: no tiene solución.
\[
\boxed{\begin{array}{ll}\lambda\neq1:&\text{solución única (SCD)},\\\lambda=1:&\text{sin solución (SI)}.\end{array}}
\]
No existe ningún valor del parámetro con infinitas soluciones.
\clearpage
\section*{Ejercicio 3(b): resolución para $\lambda=0$}
Como $0\neq1$, el caso anterior garantiza solución única. El sistema queda
\[
\begin{cases}x+y+z=2,\\2x+3z=5,\\3x+2y+4z=0.\end{cases}
\]
Restando dos veces la primera ecuación a la segunda, y tres veces la primera a la tercera, obtenemos
\[
\begin{cases}-2y+z=1,\\-y+z=-6.\end{cases}
\]
Restamos la primera de estas dos ecuaciones a la segunda:
\[y=-7.\]
Entonces
\[z=1+2y=1-14=-13,\qquad x=2-y-z=2+7+13=22.\]
\[\boxed{(x,y,z)=(22,-7,-13)}.\]
\textbf{Comprobación en las tres ecuaciones originales con $\lambda=0$:}
\begin{align*}
22-7-13&=2,\\2\cdot22+3(-13)&=44-39=5,\\3\cdot22+2(-7)+4(-13)&=66-14-52=0.
\end{align*}
Las tres igualdades se cumplen. Los valores negativos no contradicen este ejercicio algebraico: aquí las incógnitas no representan cantidades de tabletas.
\end{document}
